% Insert after the HC3 section and before the general discussion.
% Add sentiment.bib to the manuscript bibliography alongside custom.bib.
\section{Granularity in review sentiment}
\label{sec:review-sentiment}

We use review sentiment as an additional test of whether the inventory and
granularity of predicted RST relations matter for a downstream task. We
compare a representation that organizes local sentiment scores using
discourse trees with the motif representation used for HC3.

\subsection{Related work}

\citet{taboada_analyzing_2004} introduced the English SFU review
collection in a polarity classification study. On the same reviews,
\citet{voll_adjective_2007} tested adjective scoring in automatically
identified nuclei, and \citet{balage_discourse_2012} learned
relation-dependent weights using sentence-level RST annotations.
\citet{brooke_crosslingual_2009} reported 76.50\% accuracy for a unigram
SVM under ten-fold cross-validation. \citet{taboada_lexicon_2011}
reported 80.25\% accuracy for SO-CAL on these reviews, which also informed
its lexicon and system development. These scores provide historical
context under different evaluation protocols.

More broadly, \citet{bhatia_document_sentiment_2015} use discourse depth
and recursive composition to combine local sentiment in English reviews.
\citet[\S6.4.2]{sidarenka_german_twitter_2019} evaluates several relation
schemes for German Twitter sentiment by retraining a parser on the
Potsdam Commentary Corpus. Our experiment applies the HC3 label
projection to fixed predicted trees and uses simple pooling over each
tree's local sentiment scores.

\subsection{Experimental setup}

We use the 400 English Epinions reviews in the SFU Review Corpus
\citep{taboada_sfu_2008}, with 25 positive and 25 negative reviews in each
of eight product categories. The target is the reviewer's recommendation.
We use the HC3 scene-splitting approach, packing contiguous text into
chunks with a 492-token XLM-R budget and using SentSplit to find sentence
boundaries in overlong paragraphs. This yields 1,012 scenes, and 271 reviews
have multiple scenes. As in HC3, we segment scenes into
elementary discourse units (EDUs) once with PCC and give all four parsers
the same 36,850 EDUs. We use released UniRST weights for PCC and NLDT and
the fine-tuned GUM/eRST and RST-DT checkpoints described above
\citep{chistova_bridging_2025}. Each inventory supplies one tree per scene.
Fine uses its native predicted labels. Coarse maps those labels after
parsing using the HC3 mappings. EDU boundaries, attachments, and
nuclearity are identical within each fine/coarse pair.

We score each EDU independently with frozen CardiffNLP Twitter RoBERTa
\citep{loureiro_timelms_2022}.\footnote{We use
\href{https://huggingface.co/Xenova/twitter-roberta-base-sentiment-latest/tree/f3ec4d0925f90c3ca7ee7814f52d6ee7cf180445}{the pinned quantized ONNX export}.}
The \emph{equal-weight} control averages the negative, neutral, and
positive probabilities across all EDUs in the review. The
\emph{nuclearity} control takes separate averages for central and
supporting EDUs. An EDU is central when its path from its scene root
contains only nucleus branches. \emph{Relation pooling} adds six features
per relation label to the nuclearity features. At each node we average
the EDU scores in the left and right child spans separately. We then
average these span vectors over all occurrences of the label across
the review's scenes. Absent labels contribute zeros. For coarse labels
we pool the same occurrences under their mapped names.

In a second comparison we use only three-node motif frequency and
weighted average depth (MF+WAD), following the HC3 feature construction
\citep{kim_threads_2024}. Each training partition supplies a fine motif
catalog. The coarse model uses its projection under the corresponding
relation map. Relative frequencies and occurrence-weighted depths use
the same retained matches in both versions. As additional controls, we
fit logistic regression to word unigram and bigram TF--IDF features and
to log review length alone.

All classifiers use L2-regularized logistic regression and the same five
outer folds stratified by category and recommendation, with 320 training
reviews and 80 test reviews per fold. Four inner folds select
$C \in \{0.1,1,10\}$ by mean macro-F1. Vocabulary, motif catalogs, and
numeric feature scaling are fitted within the applicable training
partition. We use a fixed 0.5 decision threshold and compute macro-F1
over the 400 pooled held-out predictions. Paired, pointwise 95\% percentile
intervals use 20,000 review resamples stratified by outer fold, product
category, and gold label. Each comparison uses the same sampled reviews
for both conditions.

\subsection{Results and interpretation}

\begin{table}[t]
  \centering
  \small
  \setlength{\tabcolsep}{3.5pt}
  \begin{tabular}{@{}lrrrrr@{}}
    \hline
    & & \multicolumn{2}{c}{Relation pooling} & \multicolumn{2}{c}{MF+WAD} \\
    Inventory & Nuc. & Fine & Coarse & Fine & Coarse \\
    \hline
    PCC      & .8225 & .7900 & .7875 & .5822 & .5996 \\
    NLDT     & .8324 & .7975 & .8000 & .6025 & .5574 \\
    GUM/eRST & .8250 & .7675 & .7800 & .6074 & .5997 \\
    RST-DT   & .8224 & .7700 & .8000 & .6000 & .5822 \\
    \hline
  \end{tabular}
  \caption{Pooled held-out macro-F1 on SFU. Nuc. uses nuclearity alone to
    group EDU sentiment. Equal-weight EDU pooling, TF--IDF, and length
    score .817, .827, and .617, respectively.}
  \label{tab:sfu-granularity}
\end{table}

Separating central from supporting EDUs yields an observed gain of
0.50--1.50 macro-F1 points over equal-weight pooling across the four
inventories (Table~\ref{tab:sfu-granularity}). All relation-pooling runs
score below their nuclearity controls. With relation pooling, coarse is
numerically ahead in three inventories, with the largest gap for RST-DT
(fine minus coarse $-3.00$ points, paired interval $[-6.24,0.25]$). With
motif features, fine is ahead in three inventories, with the largest gap
for NLDT ($+4.51$ points, $[-0.23,9.27]$). Every within-inventory
fine/coarse interval and every between-inventory interval for MF+WAD
includes zero. Used on their own, motifs score below the length control.

Fine motif features yielded higher human-class F1 in every HC3 authorship
comparison (Table~\ref{tab:granularity}). On SFU, the direction of the
score difference varies with the inventory and feature construction.
The small observed nuclearity gains suggest a role for separating central
from supporting sentiment, while the fine/coarse results leave the
advantage of either granularity uncertain for these reviews.

\subsection{Limitations}

The sample covers 400 English product and service reviews, and the folds
test held-out reviews from the same eight domains. We retain the review
text as released, including any explicit ratings. The frozen scorer was
trained for Twitter and reads EDUs in isolation. Scene boundaries also
prevent discourse attachments across cuts. Between-inventory comparisons
reflect the available parser checkpoints and their training data as well
as their annotation schemes. Within-inventory coarsening holds the parses
fixed and changes feature dimensionality and sparsity alongside label
detail. Zero-filled relation columns can also carry relation-presence
information. The resampling intervals quantify variation over the saved
held-out predictions with the fitted models and split assignments held
fixed.
